\subsection{A finite weighted resolution and a proof of independence}
\label{ir:sec:resolution}

\subsubsection{The complex}
Order the primes dividing $q$. For $S\subseteq P(q)$ put
$p_S=\prod_{p\in S}p$, and define
\begin{equation}\label{ir:eq:chain-groups}
 C_j=\bigoplus_{\substack{S\subseteq P(q)\\|S|=j}}
       R_q[X_{q/p_S}].
\end{equation}
A generator is written $[S,x]$. Define
\begin{align}
 d[S,x]=\sum_{p\in S}(-1)^{\#\{\ell\in S:\ell<p\}}
 \left(\sum_{py=x}[S\setminus\{p\},y]
       -A_p[S\setminus\{p\},x]\right).
 \label{ir:eq:chain-differential}
\end{align}
The roots in each inner sum belong to $X_{q/p_{S\setminus\{p\}}}$.
The augmentation $C_0\to\mathcal D_q$ is the defining quotient.

\begin{lemma}\label{ir:lem:complex}
The operators in distinct prime directions commute before the exterior
signs are applied. Hence $d^2=0$, and the image of $C_1\to C_0$ is
exactly $\mathcal R_q$.
\end{lemma}
\begin{proof}
Both successive root sums in prime directions $p$ and $\ell$ enumerate
once all solutions of $p\ell z=x$. The two terms consisting of one
root sum and one multiplication by a weight agree in either order,
and so do the terms with coefficient $A_pA_\ell$. All intermediate
points lie in the required grids. The exterior signs make each pair
cancel in $d^2$. The assertion about $C_1$ is its definition.
\end{proof}

\subsubsection{A filtration that removes the weights}
Associate to a generator the integer
\begin{equation}\label{ir:eq:augmented-conductor}
 \nu(S,x)=p_S\operatorname{den}(x),
\end{equation}
which divides $q$. Filter the complex by increasing $\nu$.
If $p\in S$, a root of $py=x$ has denominator either
$\operatorname{den}(x)$ or $p\operatorname{den}(x)$. Thus the differential does not increase
$\nu$. Its top-filtration part retains exactly the roots of denominator
$p\operatorname{den}(x)$; the term with coefficient $A_p$ lies strictly lower.
The associated graded differential is consequently independent of all
weights.

Fix a divisor $m\mid q$. In its associated graded summand, a generator
with subset $S$ has exact denominator $m/p_S$, so necessarily
$S\subseteq P(m)$. For $p^a\parallel m$, the local coordinate belongs
to $U_{p^{a-1}}$ when $p\in S$, and to $U_{p^a}$ otherwise. At exponent
zero the unique local coordinate is zero, represented by the one-element
set $U_1$.

There is a harmless CRT twist: division by a prime changes the other
primary coordinates by a unit. We spell out its removal because omitting
it would leave a gap in a claimed tensor-product argument. If $x_\ell$
is the $\ell$-primary component of $x$, replace it in the $S$-summand by
\begin{equation}\label{ir:eq:gauge}
 z_\ell=\left(\prod_{p\in S\setminus\{\ell\}}p\right)^{-1}x_\ell.
\end{equation}
The inverse exists on the $\ell$-primary group. Under a differential
edge removing $p$, the coordinates $z_\ell$ for $\ell\ne p$ are
unchanged: the factor $p^{-1}$ in a root coordinate cancels the removal
of $p$ from the product. The $p$-coordinate is carried to all its lifts
under multiplication by $p$.

It follows that the associated graded summand at $m$ is the tensor
product, with its usual exterior signs, of the two-term complexes
\begin{equation}\label{ir:eq:local-complex}
 R_q[U_{p^{a-1}}]\xrightarrow{\,N_{p,a}\,}R_q[U_{p^a}],
 \qquad
 N_{p,a}(b)=\sum_{\substack{u\in U_{p^a}\\u\mapsto b}}u,
 \quad p^a\parallel m.
\end{equation}
For $a=1$, the image of the single generator is the sum of all $p-1$
nonzero residues. For $a\ge2$, each fiber has $p$ elements.

\begin{lemma}[Local splitting]\label{ir:lem:local-splitting}
Each $N_{p,a}$ is a split injection. Its cokernel is free with basis the
images of $E_{p,a}$.
\end{lemma}
\begin{proof}
Different fibers have disjoint supports. In one fiber, the sum of its
basis vectors has a coefficient equal to $1$ on the distinguished lift.
Replace that one vector by the fiber sum; this is a unimodular basis
change. The remaining vectors are exactly the allowed lifts. Doing this
in all fibers proves both claims over any coefficient ring.
\end{proof}

\begin{theorem}[Exact weighted resolution]\label{ir:thm:resolution}
The augmented complex
\[
 0\longrightarrow C_{|P(q)|}\longrightarrow\cdots
 \longrightarrow C_1\longrightarrow C_0\longrightarrow\mathcal D_q
 \longrightarrow0
\]
is exact. The images of $\mathcal B_q$ are a basis of $\mathcal D_q$, proving
Theorem~\ref{ir:thm:integral}.
\end{theorem}
\begin{proof}
By Lemma~\ref{ir:lem:local-splitting}, each local two-term complex is
a direct sum of a contractible identity complex and its free cokernel
in degree zero. Its finite tensor product is therefore acyclic in
positive degrees. Its degree-zero homology has the tensor-product
basis of the allowed local residues. In the graded summand at $m$,
these are exactly the basis points of denominator $m$.

Now lift through the finite filtration. If a positive-degree cycle
has highest nonzero filtration $m$, its highest component is a
boundary in the associated graded complex. Subtract a lift of that
boundary to obtain a cycle of lower filtration. Repetition makes the
cycle a boundary in the original complex. This proves positive-degree
exactness without any convergence issue.

For degree zero, suppose a linear combination of allowed points is
$d w$ with $w\in C_1$, and let $M$ be the highest filtration present
in that linear combination. If $w$ has higher filtration, its leading
component is a cycle in degree one of the associated graded complex.
Positive-degree exactness there allows us to subtract a lifted boundary
$d v$, $v\in C_2$, from $w$ and lower its filtration without changing
$d w$. Repetition gives a preimage of filtration at most $M$. Its
leading boundary would now be a relation among the allowed points in
the graded summand at $M$. Those points form a free basis of its
zero-degree homology, so all their coefficients must vanish. Descending
$M$ proves independence. Spanning was proved by
\eqref{ir:eq:straighten}. Thus the claimed points form a basis of
$H_0(C_\bullet)=\mathcal D_q$.
\end{proof}

\subsubsection{Base change and the rank identity}
The proof works verbatim over an arbitrary commutative ring with
arbitrary weights. Alternatively, the augmented complex over $R_q$
is split exact as an underlying $R_q$-module complex: its degree-zero
homology is free, and successive surjections onto projective kernels
split. Tensoring with any $R_q$-algebra therefore preserves exactness.
No flatness hypothesis on the target algebra is needed.

Taking the alternating ranks of the chain groups gives a useful check:
\begin{equation}\label{ir:eq:euler-rank}
 \rank\mathcal D_q
 =\sum_{S\subseteq P(q)}(-1)^{|S|}\frac{q}{p_S}
 =q\prod_{p\mid q}(1-p^{-1})=\varphi(q).
\end{equation}
This rank computation alone would not prove integral freeness. The
local unimodular splittings and the lifted basis are the missing
integral information.

\begin{remark}[Relation to character normal forms]
Over a characteristic-zero splitting field, both the conductor-character
basis in the source and the present raw-point basis describe the same
free module. Their transition matrix is a legitimate field-valued
change of coordinates. Its construction can involve divisions by
character-group orders. The raw-point basis is designed to avoid those
divisions, not to replace the arithmetic convenience of character
coordinates where those coordinates are available.
\end{remark}
